% !TeX root = ../papers/nc-p3.tex
\section{基底を固定した滑らかな三交換子の正確な全ファイバー}
\label{sec:marked-fixed-commutator-adjugate}
\begin{theorem}\label{thm:marked-fixed-commutator-adjugate}
標数零の代数閉体上、順序付き跡零三行列 $C_0,C_1,C_2$ を与え、
$\det\sum v_i C_i$ が滑らかな平面三次曲線を定めるとする。
$[Q_1,Q_2]=C_0$, $[Q_2,Q_0]=C_1$, $[Q_0,Q_1]=C_2$ の順序付き跡零解は
必ず存在する。基底も符号も同一視しない解集合は次の三通りで尽くされる。
\begin{enumerate}
\item 一般の場合は相異なる八解である。
\item 一つの余因子条件が階数2となる場合は相異なる六解である。
\item その階数が1となる場合は相異なる六個の孤立解と、これらと交わらない
滑らかなアフィン二次曲面 $uw-v^2=1$ が全解である。
\end{enumerate}
三通りとも実際に起こる。以下の四つの三次行列の階数で区別でき、解の復元は
楕円曲線の二等分、線型演算およびスカラー平方根で行える。
特に三交換子写像 $\mathfrak{sl}_3^3\to\mathfrak{sl}_3^3$ の一般次数は8である。
\end{theorem}
\begin{proof}[六線型座標における余因子]
前節の四つの $d$ を一つずつ扱う。Moore座標の $P=(a:b:c)$ を取り、
$D=(a(b^3-c^3),b(c^3-a^3),c(a^3-b^3))$、$M(v)=J A_D(v)$ とする。
$J$ は行 $1,2$ を交換する置換（零始まり）である。
六次元線型解を対称三行列 $Z$ で座標化する。
前節の疎な三列式で $s=(a/b,c^2/b^2,c/a)$ とした三列を $U_0$ とする
（その対角行列は零）。同式の $h=0$ を $V_0$ とし、
$W_0=-ac U_0/b^2-V_0/(3ab^4)$ と置く。$U_i,W_i$ は、三つの生成添字と
$U$ の三座標を同時に $i$ だけ巡回したものとする。
\[
            Q(Z)=\sum_{i=0}^2\big(Z_{ii}U_i+Z_{i+1,i+2}W_i\big)
\]
は線型同型である。実際各二列は対角行列を持ち得る添字ごとに三つの異なる文字部分に属し、
そこで $U_i$ は対角部分が零、$W_i$ は三乗が相異なるため非零である。

$\Lambda=\{B:\operatorname{tr}(BM(v))=0\ \forall v\}$ とする。線型同型
$\Psi:\mathrm{Sym}_3\to\Lambda$ があって
\[
        C_{Q(Z)}(v)=\Psi(\operatorname{adj} Z) M(v),\qquad
        \Psi(p p^t)=\kappa\,\operatorname{adj}(J A_P(p))
        \quad(\kappa\ne0)
\]
となる。次の不定元計算で恒等式と可逆性を確認する。
$M(e_0)$ は可逆なので左倍率を $C_{Q(Z)}(e_0)M(e_0)^{-1}$ と定める。
$Z$ の二次単項式21個の係数ごとに、他の二スライスが同じ左倍率である。
$Z=p p^t$ では四次単項式15個の係数は全て零になる。対称三次行列の二次式のうち
$pp^t$ で消えるものは六つの余因子で張られる。実際二次単項式21個から
四次15個への代入は全射で、六余因子は独立である。
$Z_{i+1,i+2}^2$ と $Z_{ii} Z_{i+1,i+2}$ の左倍率係数を取れば、それぞれ符号を除き
$\Psi(E_{ii})$, $\Psi(E_{i+1,i+2}+E_{i+2,i+1})$ となる。
$i=0,1,2$ でこの二列の支持は互いに素な三本の巡回対角であり、
各対の最初の二成分の小行列式は順に
\[
       (a^3-c^3)/(a^4b^7),\quad -c(a^3-b^3)/(a^4b^8),
                      \quad -(b^3-c^3)/(a^3b^8)
\]
である。よって六列は独立、像は跡条件の六空間に等しい。
同じ六列を $p_ip_j$ で足すと上記のMoore行列の余因子になる。
\path{computations/verify_fixed_commutator_adjugate.py}はこれらを全係数の有理恒等式として検算する。

入力を $C(v)=L M(v)R$ と書く。前節の線型比較で可逆な $L,R$ は
$(L,R)\mapsto(\lambda L,\lambda^{-1}R)$ を除き一意である。
この補助類の実際の解は
\[
       S=\Psi^{-1}(RL),\qquad \operatorname{adj} Z=S,\qquad
                      Q_i^{\rm actual}=R^{-1} Q_i(Z) R
\]
で全て得られる。跡零条件により $RL\in\Lambda$ である。
十分性は恒等式による。必要性について、補助表示の行基底は $W$ を変えず、
右基底の変更は $Q$ の同時共役により任意に動かせる。
上の積表示から全ての非特異解の右核点は座標ごと $M$ の右核点であるため、
入力に戻すその射影的な共役変換は $R^{-1}$ に限られる。
実際右核点は平面三次曲線を埋め込み、その全点を固定する射影変換は恒等である。
ゆえに独立な左右基底を誤って許すことなく、表示した全解となる。
$S$ は非零で、階数3なら $Z=\pm\sqrt{\det S}\,S^{-1}$ の二解、
階数2なら解なし、階数1なら合同座標で非退化な二次対称ブロックの
行列式を固定した解全体である。後者では $ZS=0$ から他の行と列が零になる。
\end{proof}

\begin{proof}[異なる二選択は同時に退化しない]
入力曲線とその\emph{具体的な右核点写像}を固定し、ある代表 $M_*$ を選んで
入力を $B M_*$ と動かす。可逆な $B$ は $\Lambda_*=\{B:\operatorname{tr}(BM_*)=0\}$ に入る。
この六空間の階数一行列 $u x^t$ は $x^t M_*(v)u=0$ を全 $v$ で満たす。
その射影軌跡 $E'$ は元の曲線 $E$ に同型である。実際 $u$ は $M_*$ の右核点の
平面三次曲線を動き、$M_*(-)u$ の階数は2なので $x$ は一意である。
後者は変数と右列の役割を替えたMoore表示でも行列式が同じ滑らかなHesse三次式に
比例することによる。両核写像は正則である。この同定による列点の線束を
$T_{\rm col}$、Segre埋込みの線束を $L_6$ とする。

四つの $\Psi$ を同じ $\Lambda_*$ に戻す操作は左右の定数行列によるので階数を保つ。
詳しくは $M_*=L_0 M R_0$ なら $B\mapsto R_0 B L_0$ である。
$\Psi(pp^t)$ の公式により、それぞれのVeronese曲面は $E'$ 全体を含む。
その曲面の平面座標から来る $E'$ 上の線束を $N_d$ とすると $2N_d=L_6$ である。
これらは相異なる。実際 $p\in E$ による $E'$ の表示で、行列の列点は
$A_P(p)$ の右核であり、元の $E$ への同定で使う列点は $A_D$ の右核である。
上の左右基底 $R_0,L_0$ はこの比較では打ち消し合う。
Hesse曲線上の $p\mapsto\ker A_P(p)$ は平行移動である。
変数と列を交換しても同じHesse曲線となり、$P$ を変曲点も含めて動かすと
自己同型の代数的な族で、$P=O$ では恒等である。楕円曲線の自己同型の恒等成分は
平行移動群なので主張が従う。$A_D$ にも同じことが成立する。
ゆえに二同定のずれも平行移動で $\mathrm{Pic}^0$ には恒等に作用し、
$T_{\rm col}-N_d=d$ である。四つの $d$ は相異なるので全四平方根を得る。
さらに $E'$ はこれらの曲面上で平面三次曲線であり、$L_6$ の完全系で埋め込まれている。

$S$ が特異である条件は、$B$ の点が対応する $N_d$ のある三次因子の
$L_6$ による包に入ることである。すなわち $B$ が $H^0(N_d)$ 上で定める
対称形式の核に非零切断 $s$ があり、$sH^0(N_d)$ を消すことと同値である。
$B$ は可逆なので、階数一行列の曲線 $E'$ の長さ2以下の部分の包には入らない。
異なる平方根 $N,N'$ の三次因子がともに $B$ を張るなら、
定理\ref{thm:marked-fixed-metric-elliptic}の割線の議論で両因子は交わらず
$N+N'=L_6$ となるが、$2N=L_6$ に反する。
従って四選択のうち特異なのは高々一つであり、定理の個数が従う。

各例外の実在については、$p\notin E$ の $pp^t$ が可逆な $\Psi(pp^t)$ を与えるので
階数1が実現する。そこから対称行列を階数2へ変形しても一般点で倍率は可逆なままである。
四つの固有な三次超曲面の外は非空開であり、一般の順序付き跡零標的では全てが階数3である。
各二解の局所的な復元はこの開集合でエタールである（標数零の二等分と非零平方根、
および単元小行列での線型解）。従って一般次数も8となる。
\end{proof}
